\section{Method}
\label{sec:method}

We target dense decoder-only LLMs and prune two kinds of structured units: attention units
(a head in multi-head attention, or a KV-aligned query-head group in grouped-query
attention~\citep{ainslie2023gqa}) and FFN groups (contiguous blocks of intermediate
channels). Let $\Uset=\{u_1,\dots,u_M\}$ collect all such units across all layers, with a removable
parameter cost $c_u=\Param(u)$ each. Each unit carries a keep gate $m_u$ scaling its write into the residual stream, and we write
$s_u=1-m_u$ for the pruning variable: $s_u{=}0$ keeps unit $u$ and $s_u{=}1$ zeroes its write. The
deployed mask is binary, $s\in\{0,1\}^M$. The gate is real-valued, so $\mathcal{R}$ below is defined on
a neighbourhood of $\mathbf 0$ in $\mathbb{R}^M$ and $\mathbf 0$ is an interior point rather than a
vertex of the cube, which is what makes its expansion well posed. Speedups are measured on models whose weight matrices are
physically sliced, not on masked full-size ones (\cref{app:units}).

\subsection{Pruning risk as second-order self-distillation}
\label{sec:method-risk}

Let $f_0$ be the frozen full model (the teacher) with next-token distribution
$p_0(\cdot\mid x_{\le t})=\mathrm{softmax}(z_0(x,t))$ over a vocabulary of size $V$, and let $f_s$ be
the masked model with distribution $p_s$. Given a small unlabeled calibration set
$\mathcal{C}$ and a deterministic set of token positions $\mathcal{T}(x)$ per sequence, we define the
pruning risk as the token-level self-distillation KL divergence
\begin{equation}
\mathcal{R}(s)=\mathbb{E}_{x\in\mathcal{C}}\Big[\tfrac{1}{|\mathcal{T}(x)|}
\textstyle\sum_{t\in\mathcal{T}(x)}\KL\big(p_0(\cdot\mid x_{\le t})\,\|\,p_s(\cdot\mid x_{\le t})\big)\Big].
\label{eq:risk}
\end{equation}
This objective is label-free and ties directly to output behaviour. Its defining property is its local
geometry. The masked model equals the teacher at $s=\mathbf{0}$, so the two distributions coincide there. Since
$\KL$ is non-negative and vanishes only at equality, $\mathbf 0$ is an interior minimiser of
$\mathcal{R}$, and both the value and the gradient vanish there:
\begin{equation}
\mathcal{R}(\mathbf{0})=0,\qquad \nabla_s\mathcal{R}(\mathbf{0})=0.
\label{eq:vanish}
\end{equation}
A second-order Taylor expansion around the full model has no first-order term, and the leading
behaviour is a pure quadratic:
\begin{equation}
\mathcal{R}(s)\;\approx\;\tfrac{1}{2}\,s^\top \Hmat^\star\, s,
\qquad
\Hmat^\star=\nabla_s^2\mathcal{R}(s)\big|_{s=\mathbf{0}}\in\mathbb{R}^{M\times M}.
\label{eq:quadratic}
\end{equation}
For binary masks $s_u^2=s_u$, so expanding \cref{eq:quadratic} gives the central decomposition
\begin{equation}
\mathcal{R}(s)\approx
\underbrace{\tfrac12\textstyle\sum_{u}\Hmat^\star_{uu}\,s_u}_{\text{node saliency (diagonal)}}
\;+\;
\underbrace{\textstyle\sum_{u<v}\Hmat^\star_{uv}\,s_u s_v}_{\text{co-pruning edges (off-diagonal)}}.
\label{eq:decomp}
\end{equation}

\begin{remark}[Node saliency is the diagonal special case]
Dropping the off-diagonal terms reduces \cref{eq:decomp} to a sum of independent saliencies
$\tfrac12\sum_u \Hmat^\star_{uu}s_u$, i.e.\ an Optimal-Brain-Damage--style criterion~\citep{lecun1989obd}
lifted to structured units (formalized in \cref{prop:obd}). The cross-module edges $\Hmat^\star_{uv}$ encode how the
removal of $u$ and $v$ interacts: $\Hmat^\star_{uv}>0$ means co-pruning is more harmful than the
units' individual terms predict; $\Hmat^\star_{uv}<0$ means their perturbations partly cancel under the
Fisher metric; $\Hmat^\star_{uv}\approx 0$ means they are second-order independent. A low-saliency
\emph{bridge} unit (small $\Hmat^\star_{uu}$) with many strong edges is thus protected by an accurate estimator, because pruning it
together with its partners is costly, exactly the failure mode node-first scoring misses.
\end{remark}

\subsection{Estimating the edge matrix from single-unit ablations}
\label{sec:method-estimate}

The exact Hessian $\Hmat^\star$ in \cref{eq:quadratic} \emph{is} the Fisher curvature of \cref{eq:risk}, not a
Gauss--Newton approximation of it: the KL gradient in logit space vanishes at $s=\mathbf 0$
(\cref{prop:vanish}), so the term a Gauss--Newton step would drop is exactly zero there. And it has a
closed form we can evaluate. For one token position the KL Fisher in logit space is
$F_{x,t}=\mathrm{diag}(p_0)-p_0 p_0^\top$, and to second order the KL between teacher and masked model is the Fisher-weighted norm of the logit
perturbation $\Delta z_s:=z_0-z_s$, that is $\KL(p_0\|p_s)\approx\tfrac12\,\Delta z_s^\top F\,\Delta z_s$
(\cref{prop:fisher}). The key practical step is that this perturbation is, to leading order,
additive across units: if $\delta z_u(x,t)=z_0(x,t)-z_{-u}(x,t)$ is the logit shift from
masking unit $u$ alone, then $\Delta z_s\approx\sum_u s_u\,\delta z_u$. Writing
$J_u=\partial z_s/\partial s_u|_{s=\mathbf 0}$, the exact entry is
$\Hmat^\star_{uv}=\mathbb E_{x,t}[J_u^\top F_{x,t}J_v]$. Replacing $J_u$ by the observed
full-ablation displacement $\delta z_u$ gives the finite-ablation estimator used by the algorithm,
\begin{equation}
\Hmat_{uv}:=\mathbb{E}_{x,t}\big[\,\delta z_u(x,t)^\top F_{x,t}\,\delta z_v(x,t)\,\big].
\label{eq:huv}
\end{equation}
Whitening the perturbation into its probability-weighted, mean-centred form
$\widetilde{\delta z}_u=\sqrt{p_0}\odot(\delta z_u-\mathbb{E}_{y\sim p_0}[\delta z_{u,y}])$ turns the
Fisher inner product into a plain dot product (\cref{app:estimation}), giving
\begin{equation}
\Hmat_{uv}=\tfrac{1}{P}\,\bar{D}_u^\top \bar{D}_v,
\qquad \bar{D}_u=\mathrm{vec}\big(\{\widetilde{\delta z}_u(x,t)\}_{x,t}\big),
\label{eq:gram}
\end{equation}
where $\bar D$ is the matrix with columns $\bar D_u$ and $P$ is the total number of calibration
positions: we keep a fixed number of positions per sequence, so $P$ is that count times the number of
sequences ($128\times128$ here) and pooling the positions is the same average as the nested one in
\cref{eq:risk}. The entire $M\times M$ edge matrix is therefore a Gram matrix of \emph{single-unit}
feature vectors. It estimates $\Hmat^\star$ with $M$ single-unit ablation passes, one per unit, rather than $O(M^2)$ pairwise
ablations, and needs no gradients or Hessian-vector products. From here on $\Hmat$ denotes this
estimated matrix, and
\cref{app:th-est} separates the two approximations involved and measures both. The $V\times V$ Fisher is never materialised: features live on the teacher's top-$r$ logit
coordinates, renormalised to sum to one and shared across units, so the Gram remains valid and $F$ is
the Fisher of that truncated distribution (\cref{rem:exact}).
\Cref{app:estimation} gives the implementation and its complexity.

\subsection{One-shot cost-normalised greedy co-pruning}
\label{sec:method-solver}

Write $\widehat{\mathcal R}_\lambda(s)=\tfrac12\sum_u\Hmat_{uu}s_u+\lambda\sum_{u<v}\Hmat_{uv}s_us_v$ for
the surrogate risk of \cref{eq:decomp}, with the interaction strength $\lambda$ of
\cref{sec:method-lambda}. Given a target ratio $\rho$, pruning solves the budgeted binary quadratic
program
\begin{equation}
\min_{s\in\{0,1\}^M}\;\widehat{\mathcal R}_\lambda(s)
\quad\text{s.t.}\quad \textstyle\sum_{u}c_u s_u\;\ge\;\rho\textstyle\sum_u c_u .
\label{eq:objective}
\end{equation}
Attention units and FFN groups share one budget pool, so how the budget divides between the two types
emerges from the optimisation rather than being fixed. Exact solution is intractable. We use a one-shot
greedy solver on the fixed surrogate, with no weight updates and no re-estimation of $\Hmat$.
Maintaining the cumulative perturbation of the current set $S$, the marginal risk of adding unit $u$
is
\begin{equation}
\Delta R(u\mid S)=\tfrac12 \Hmat_{uu}+\textstyle\sum_{v\in S}\Hmat_{uv}=\tfrac12\Hmat_{uu}+g_u^{(S)},
\label{eq:marginal}
\end{equation}
where $g_u^{(S)}=\sum_{v\in S}\Hmat_{uv}$ is a running interaction accumulator updated in $O(M)$ by
$g\!\leftarrow\! g+\Hmat_{:,u^\star}$ after each selection. At each step we pick the unit with the
smallest cost-normalised marginal risk, $u^\star=\arg\min_{u\notin S}\Delta R(u\mid S)/c_u$, and
stop once the budget is met. The solver is $O(M|S|)$ for $|S|$ selected units, and \cref{alg:edge-taylor} (\cref{app:solver})
summarises the full pipeline. A per-layer cap keeps removal from concentrating; it is shared by every
method we compare, at every ratio, and specified in \cref{app:protocol}.

\subsection{Interaction strength: a path, not a hyperparameter}
\label{sec:method-lambda}

The diagonal follows from single-unit ablations alone. The off-diagonal additionally assumes that
simultaneous removals perturb the logits additively (\cref{eq:huv}), so we weight it by a single
\emph{interaction strength} $\lambda\in[0,1]$,
\begin{equation}
\widehat{\mathcal R}_\lambda(s)=\tfrac12\textstyle\sum_u\Hmat_{uu}s_u
+\lambda\textstyle\sum_{u<v}\Hmat_{uv}s_u s_v,
\label{eq:lambda}
\end{equation}
so that $\lambda\!=\!0$ recovers pure node saliency and $\lambda\!=\!1$ the full edge model. The $\lambda\!=\!0$ endpoint is a floor rather than an arbitrary limit. It is exactly the OBD diagonal
selector (\cref{app:classical}), so whatever $\lambda$ we end up with, the method degrades at worst to
a strong classical second-order baseline.

Choosing it is not a fitting problem. The solver consumes only the \emph{order} the marginals
induce, so $\lambda$ has to be chosen by what the selections it produces actually cost --- and that
quantity turns out to be exactly computable.

\textbf{The family of solutions is finite.} What the solver does consume is the ordering of the
conditional marginals $\Delta_\lambda(u\mid S)=\tfrac12\Hmat_{uu}+\lambda\,g_u^{(S)}$. Normalised by
cost, each candidate contributes an affine function of $\lambda$,
\begin{equation}
\Delta_\lambda(u\mid S)/c_u \;=\; \mu_u+\lambda\,\nu_u,
\qquad \mu_u=\tfrac{\Hmat_{uu}}{2c_u},\qquad \nu_u=\tfrac{g_u^{(S)}}{c_u},
\label{eq:affine}
\end{equation}
so within a single greedy step the arg-min can change only where two of these lines cross. Every step of the solver therefore switches at finitely many values of $\lambda$, and the solution
$S_\lambda$ is \emph{piecewise constant}. The family $\{S_\lambda:\lambda\in[0,1]\}$ is a finite set of
prunings, a few hundred per model rather than a continuum.

\textbf{The family can be enumerated exactly.} Because \cref{eq:affine} is affine, the crossings are
available in closed form. Starting at $\lambda\!=\!0$, the earliest crossing over all greedy steps is the right endpoint of the
interval on which the current solution is optimal. Stepping just past it and re-solving moves to the
next member, and repeating walks the whole family; the step past each crossing is $10^{-9}$, nearly
three orders of magnitude below the narrowest interval on which any member of the ten enumerated
families is optimal, $7.6\times10^{-7}$ (\cref{app:lambda}). The construction is the same critical-value walk that makes the
LASSO path enumerable by LARS~\citep{efron2004lars}, borrowed for its mechanics rather than its
guarantees: LARS traces the exact optimum of a convex problem and its path is piecewise linear, while
$S_\lambda$ is the path of a greedy heuristic and is piecewise constant. The walk is CPU-only, one greedy solve per breakpoint and no
forward passes, and what the selection then costs is one masked forward pass per distinct
\emph{member} --- $303$ of them on Llama-3.1-8B-Instruct. Sampling $\lambda$ is not a substitute at any
price. The median interval width is $1.4\times10^{-3}$ across the ten models, so a grid at step $0.002$
still reaches a median of only $61\%$ of the members and as few as $35\%$: it misses solutions rather
than costing more (\cref{app:lambda}).

\textbf{Selection is then a held-out measurement.} With the family in hand there is nothing left to
search over; we keep the member with the lowest measured held-out risk,
\begin{equation}
\hat\lambda=\arg\min_{\lambda\in[0,1]}\ \mathcal{R}(S_\lambda;\mathcal{C}'),
\label{eq:lambdahat}
\end{equation}
where $\mathcal{R}(\cdot\,;\mathcal{C}')$ is the same token-level KL of \cref{eq:risk}, measured on a
calibration split $\mathcal{C}'$ held out from the one that built $\Hmat$. No labels and no benchmark enter, and each $S_\lambda$ is judged by measured risk rather than by the
surrogate that proposed it. Every member is scored on the same sequences, so differences across
$\lambda$ are paired and we report them in paired standard errors. The cost is one masked forward pass
per member over a small held-out set, against the $M$ ablations $\Hmat$ already requires.

Two properties of this rule matter for how the results below should be read. First, the setting is
fixed before any benchmark is seen: $\hat\lambda$ is chosen on held-out risk, so no test set enters the
selection, and we deliberately do not pick the member that scores best downstream --- a family member
tuned on perplexity would be a different and much weaker claim. Second, the criterion is flat near its
optimum. Neighbouring members sit within its own noise: on the headline model the member at
$\lambda{=}0.15$ carries $2.1\%$ more held-out risk, $1.3$ paired standard errors, and lands at a
comparable operating point. The rule therefore does not need to be delicate to be useful, and
\cref{app:lambda-robust} confirms that resampling the held-out set and changing the selection protocol
move it very little.

\textbf{What the resulting curve says.} The diagonal contributes one term to the marginal while the
off-diagonal contributes $|S|$ of them, so a balanced combination would sit near
$\lambda\sim 1/(2|S|\bar\rho_S)$, with $\bar\rho_S$ the mean curvature correlation over in-set pairs
(\cref{prop:offdiag}) --- of order $10^{-2}$ at the ratios we prune. The criterion lands at or
below that scale: $\hat\lambda<0.1$ on nine of the ten models and exactly $0$ on two. The off-diagonal
is worth consulting and is damped hard, and at $\hat\lambda$ it still supplies a median $12\%$ to
$46\%$ of the winning score. $\mathcal{R}(S_\lambda)$ has an \emph{interior} minimum on eight of the ten models we measure. Both $\lambda\!=\!0$ (no edges) and $\lambda\!=\!1$ (full trust) are
worse than the calibrated combination, which is direct evidence that the off-diagonal carries signal
and that it must nonetheless be damped. On the other two the selected segment begins at
$\lambda\!=\!0$ and the procedure returns the diagonal selector unchanged. That is what \cref{eq:lambdahat} buys: the
diagonal solution is itself a member of the family being minimised over, so the criterion cannot
return something with higher held-out risk than node saliency alone. We measure $\mathcal{C}'$ at the deployment
context length, so the risk we minimise is the risk the pruned model will incur; the selection is
insensitive to that choice either way (\cref{app:lambda}).

\paragraph{One matrix, two uses.} The couplings that say what not to remove together also say how the
survivors should absorb what was removed. The same blocks of $\Hmat$ give a coupling-aware
compensation on the kept FFN units in closed form, from one small ridge solve (\cref{app:comp}). Every
number outside \cref{sec:exp-recovery} is selection-only.
